\runninghead{Scope and References}
\section{Appendix: Scope and Qualifications}
\label{sec:scope}

{\small
\textbf{Edition, formulary and occurrence.} The study treats the Mass
\latin{D.~N. Iesu Christi Regis}, \latin{Dominica ultima octobris}, I classis,
of the \work{Missale Romanum}, editio typica, Typis Polyglottis Vaticanis 1962,
printed pp.~697--698, nos.~3938--3949, with its proper Preface, p.~291,
no.~1080, in the universal 1962 calendar and no particular calendar. Its
occurrence on Sunday 25~October 2026, its displacement of the Twenty-second
Sunday after Pentecost without commemoration, the omission of the
commemoration of Saints Chrysanthus and Daria, the colour, the Gloria, the
Credo and the Preface rest on the Missal's own calendar and general rubrics
(RG 17~d, 91, 111~a, 112~b, 120~a; RGMR 431, 475, 492), as
\texttt{research/context.md} records. The Tract and the paschal Alleluia
printed with the formulary belong to votive use on other days and are given as
printed text, not as parts of this Mass. The \latin{Asperges}, the Ordinary,
and the act of consecration of mankind to the Sacred Heart that \work{Quas
primas} n.~28 prescribes for the day are outside the formulary and are not
treated.

\textbf{Text control.} Every Latin form was read on page images of the CMAA
facsimile of the 1962 typical edition and collated word for word with the
formulary's first publication in \work{Acta Apostolicae Sedis} 17 (1925),
pp.~666--668; the collation, its three bounded exceptions (the conclusions,
the Preface's common closing, the last word of the paschal Alleluia) and the
comparison with the Clementine Vulgate are in \texttt{propers/verified.md}.
The 1960 rubrics printed with the formulary are described and not reproduced.
The facsimile's page images and its engraved headpiece are not reproduced.

\textbf{English.} The scriptural English is the Douay--Rheims in Challoner's
revision at the canonical verses, with the Missal's departures named beside
each text; the Challoner text numbers the Introit's verse Ps 71:2 and prints
the Communion's two half-verses inside one verse numbered 28:10. The English
of the orations and the Preface is the Lasance \work{New Roman Missal} of 1945,
whose rubrics are those before 1960. Neither is an approved liturgical
translation of the 1962 Missal. Patristic and later witnesses are quoted in
the nineteenth-century English of the editions named in the References, or in
their Latin, and otherwise described; no English is composed for a liturgical
text.

\textbf{Reception.} The readings rest on the loci named in the text and
recorded, with editions, digests and bounds, in \texttt{research/scope.md}
and \texttt{research/interpretations.md}. The sweep is bounded by the corpora
named there and is not complete. The Collect, Secret, Preface and
Postcommunion were composed for the feast in 1925, and no patristic witness
comments on them as texts; their verbal sources are named where they bear on a
claim. Blessed Ildefonso Schuster's \work{The Sacramentary} V (English, 1930)
is the one commentary read on this Mass as a Mass; he comments on its elements
one by one, and the one join between two elements he makes himself is the
Gradual with the Alleluia. Every other connection drawn between elements of
the formulary belongs to the study's readings unless a witness is named as
making it. Bellarmine is read in O'Sullivan's abridged English of 1866,
checked against its facsimile; Cassiodorus and Aquinas in modern web
transcriptions of Migne's, Parma's and Marietti's texts; Theodoret in the PG
printings, partly on their text layers; Augustine's \work{Letter}~185 and
Chrysostom's homilies on John in web transcriptions of the NPNF English.
Cassiodorus and Theodoret of Cyrus are ecclesiastical writers, and St Cyril of
Alexandria, St Cyril of Jerusalem, the Daniel scholia under Hippolytus's name
and Victorinus of Pettau support particular claims; none of them carries a
reading. \work{Quas primas} is cited as the instituting act of the feast, by
number, in the Holy See's English.

\textbf{Chronology.} Every biblical date in the Scriptural Date and Location
appendix is generated from the repository's Scripture chronology corpus,
with its sources named there. The corpus reads the formulary from the 1962
calendar finding aid, which stores the votive Tract as an ordinary chant of
the day and does not carry the paschal Alleluia; the Tract's rows are printed
with the rubric that limits it, and Apoc 19:16 has no row. The record holds no
date for the questioning before Pilate or for the first year of Baltasar.

\textbf{Rights.} The Latin is published on 17 U.S.C. 103(b) together with its
first publication abroad in 1925, before the 1931 line recorded in
\texttt{guidance/missals.md}, applied as of 19~October 2026. The Douay--Rheims
and the Lasance English (public domain in the United States by nonrenewal)
are quoted from registered witnesses; the NPNF, ANF, LFC, Migne and CSEL
editions are nineteenth-century publications, and Schuster's volume was first
published in London in 1930. Quotations from the Corpus Thomisticum
transcriptions and the Holy See's English of \work{Quas primas} are kept brief.
Scripture, liturgical texts and quoted translations are not relicensed by the
repository's licence.

\textbf{Use.} This is a study aid to the 1962 Missal, not an official
liturgical book, translation or critical edition.
\par}

\section{References}
\label{sec:references}

{\small
\begin{itemize}[leftmargin=1.2em,itemsep=.25em]
\item \work{Missale Romanum ex decreto Sacrosancti Concilii Tridentini
restitutum, Summorum Pontificum cura recognitum}, editio typica (Typis
Polyglottis Vaticanis, 1962), CMAA facsimile: \work{Rubricae generales} and
\work{Rubricae generales Missalis romani}; Sanctorale, pp.~697--698,
nos.~3938--3949; \work{Praefationes}, p.~291, no.~1080.
\item \work{Acta Apostolicae Sedis} 17 (1925), fasc.~15, pp.~666--668: Sacred
Congregation of Rites, decree \latin{Urbis et Orbis} of 12~December 1925,
with the Mass of the feast; Internet Archive microfilm scan.
\item \work{Missale Romanum} (Ratisbon: Pustet, 1862), the common closing of
the Preface.
\item Lasance, F.~X., \work{The New Roman Missal in Latin and English}, revised
edition (New York: Benziger Brothers, 1945), pp.~813--814 (Preface of Jesus
Christ, King) and 1257--1260 (Feast of Our Lord Jesus Christ, King).
\item \work{The Holy Bible Translated from the Latin Vulgate} (Douay--Rheims,
Challoner revision), Project Gutenberg text: Ps 2:6--11; 28:1--10; 44:8;
71:2, 8--11, 20; 88:1, 27--30, 39--45; 109:3; Wis 8:1; Dan 1:1--6; 7:1,
13--14; Mt 28:18; Lk 1:33; Jn 17:1--2, 20; 18:33--37; Acts 13:33; 1 Cor
15:24--28; 2 Cor 5:18; Phil 2:6--8; Col 1:12--20; 2 Tim 2:12; Heb 2:8; Apoc
1:5--6, 9; 5:5--12; 19:16.
\item \work{Biblia Sacra Vulgatae Editionis} (Clementine Vulgate), eBible
text, for the comparisons with the Missal's forms.
\item Pius~XI, \work{Quas primas}, encyclical of 11~December 1925, English on
the Holy See's website, nn.~13, 16--19, 28, 32.
\item Augustine, \work{Tractates on the Gospel of John} 115.1--4, NPNF, first
series, vol.~7 (CCEL text).
\item Augustine, \work{Enarrationes in Psalmos} (\work{Expositions on the Book
of Psalms}), NPNF, first series, vol.~8 (CCEL text): 2.6--9; 28.10--11; 71.1,
11--14.
\item Augustine, \work{Letter} 185, \work{The Correction of the Donatists},
sections 19--20, trans. J.~R. King, NPNF, first series, vol.~4 (New Advent
text).
\item Augustine, \work{De Trinitate} I.8.15--17, I.10.20--21, trans.
A.~W. Haddan (New Advent text).
\item Basil the Great, \work{Homilia in Psalmum XXVIII}, section 8, PG 29,
cols.~304C--305A (Internet Archive page images).
\item Bede, \work{Explanatio Apocalypsis}, book I, on 1:5--6 and 5:5--7, PL 93
(Internet Archive OCR), with Marshall's English of 1878.
\item Bellarmine, Robert, \work{Explanatio in Psalmos}, in John O'Sullivan's
English, \work{A Commentary on the Book of Psalms} (Dublin: James Duffy, 1866),
on Pss 2, 28 and 71.
\item Cassiodorus, \work{Expositio Psalmorum}, on the second, twenty-eighth and
seventy-first psalms, Migne's text (Corpus Corporum transcription).
\item Chrysostom, John, \work{Homilies on John}, homilies 83--84, NPNF,
first series, vol.~14 (New Advent text).
\item Chrysostom, John, \work{Homilies on Colossians}, homilies 2--3, NPNF,
first series, vol.~13 (CCEL text).
\item Cyril of Alexandria, \work{Commentary on John}, book XII, on 18:36--37,
Library of the Fathers 48, pp.~599--601 (tertullian.org transcription).
\item Cyril of Jerusalem, \work{Catechetical Lectures} 15.27--32, NPNF, second
series, vol.~7 (CCEL text), with the editor's note identifying the teaching
opposed.
\item Hilary of Poitiers, \work{Tractatus super Psalmos} 2.30--33, 2.42--43,
ed. A.~Zingerle, CSEL 22 (1891), pp.~59--62, 70 (Internet Archive page images).
\item Hilary of Poitiers, \work{De Trinitate} XI.39--40, NPNF, second series,
vol.~9 (CCEL text).
\item Hippolytus (scholia transmitted under his name), on Daniel, ANF, vol.~5
(CCEL text).
\item Jerome, \work{Commentarii in Danielem}, book II, on chapter seven, PL 25,
cols.~670--671 (Wikisource transcription of Migne).
\item Schuster, Ildefonso, \work{The Sacramentary (Liber Sacramentorum)},
vol.~V (London: Burns Oates \& Washbourne, 1930), pp.~200--204 (Internet
Archive scan).
\item Theodoret of Cyrus, \work{Interpretatio in Psalmos}, on the second,
twenty-eighth and seventy-first psalms, PG 80; \work{Interpretatio Epistolae
ad Colossenses}, chapter one, PG 82.
\item Thomas Aquinas, \work{Super Evangelium S.~Ioannis lectura}, chapter 18,
lect.~6; \work{Postilla super Psalmos}, on the second and twenty-eighth
psalms; \work{Super Epistolam ad Colossenses lectura}, chapter one,
lect.~3--5 (Corpus Thomisticum transcriptions).
\item Victorinus of Pettau, \work{Commentary on the Apocalypse}, on chapter
five, ANF, vol.~7 (New Advent text).
\item \work{The Catholic Encyclopedia} (New York: Robert Appleton, 1907--1912):
Aherne, ``Epistle to the Colossians'' (vol.~4); Durand, ``The New Testament''
(vol.~14); Fonck, ``Gospel of St.~John'' (vol.~8); Gigot, ``Book of Daniel''
(vol.~4); Meistermann, ``Temple of Jerusalem'' (vol.~14); Prat, ``St.~Paul''
(vol.~11); Reid, ``Captivities of the Israelites'' (vol.~3); Schets, ``Third
and Fourth Books of Kings'' (vol.~8); Sloet, ``Chronology of the Kings''
(vol.~8); Van den Biesen, ``Apocalypse'' (vol.~1).
\item Haydock, George Leo, \work{The Holy Bible} with Haydock's commentary
(Loreto reprint, 2014), the chronological paragraph attributed to Ussher in its
commentary on the seventieth psalm.
\item \work{New American Bible, Revised Edition}, introduction to the Psalms
(USCCB website), summarized.
\end{itemize}
\par}
